\documentclass[12pt]{article}
\usepackage{fullpage}
\usepackage{amsmath,amssymb,amsthm}
\usepackage{hyperref}
\usepackage{amsfonts}

\newtheorem{theorem}{Theorem}
\newtheorem{proposition}[theorem]{Proposition}
\newtheorem{lemma}[theorem]{Lemma}
\newtheorem{corollary}[theorem]{Corollary}
\theoremstyle{definition}
\newtheorem{definition}[theorem]{Definition}
\newtheorem{remark}[theorem]{Remark}
\newtheorem{example}[theorem]{Example}

\DeclareMathOperator{\SO}{SO}
\DeclareMathOperator{\SL}{SL}
\DeclareMathOperator{\GL}{GL}
\DeclareMathOperator{\Hom}{Hom}
\DeclareMathOperator{\vcd}{vcd}
\DeclareMathOperator{\cd}{cd}

\title{Problem 7: Uniform Lattices with $2$-Torsion as Fundamental Groups\\
of Manifolds with $\mathbb{Q}$-Acyclic Universal Cover}
\author{}
\date{}

\begin{document}
\maketitle

\section*{Statement and Answer}

\noindent\textbf{Question.}
Suppose that $\Gamma$ is a uniform lattice in a real semisimple group, and that $\Gamma$
contains some $2$-torsion.  Is it possible for $\Gamma$ to be the fundamental group of
a compact manifold without boundary whose universal cover is acyclic over $\mathbb{Q}$?

\medskip
\noindent\textbf{Answer.}  \textbf{Yes.}

\section{Setup and Definitions}

We fix notation throughout.

\begin{definition}
A \emph{uniform lattice} in a locally compact group $G$ is a discrete subgroup
$\Gamma \leq G$ such that the quotient $\Gamma\backslash G$ is compact.  A real semisimple Lie group
$G$ is a connected Lie group whose Lie algebra $\mathfrak{g}$ is semisimple, with no compact
normal subgroups of positive dimension.
\end{definition}

\begin{definition}
Let $M$ be a connected manifold with fundamental group $\pi_1(M) \cong \Gamma$ and
universal cover $\widetilde{M}$.  We say $\widetilde{M}$ is \emph{$\mathbb{Q}$-acyclic} if
\[
  H_k(\widetilde{M};\mathbb{Q}) \;=\; 0 \quad \text{for all } k \geq 1.
\]
\end{definition}

\begin{definition}
A group $\Gamma$ has \emph{$2$-torsion} if there exists $\gamma \in \Gamma$
with $\gamma \neq e$ and $\gamma^{2} = e$.
\end{definition}

\begin{definition}
Let $\Gamma$ be a group of type $\mathrm{FP}$ (i.e., admitting a projective resolution of
$\mathbb{Z}$ over $\mathbb{Z}\Gamma$ that is finitely generated in each degree).  We call
$\Gamma$ an \emph{$n$-dimensional duality group} if
\[
  H^{i}(\Gamma;\,\mathbb{Z}\Gamma) \;=\;
  \begin{cases} D & i = n, \\ 0 & i \neq n,\end{cases}
\]
for some $\mathbb{Z}\Gamma$-module $D$ called the \emph{dualizing module}.
\end{definition}

\begin{definition}
A duality group is of \emph{finite type} if it is finitely presented and has a
$K(\Gamma,1)$-complex with finitely many cells in each dimension (i.e., $\Gamma$ is of
type $F_\infty$).
\end{definition}

\section{Construction: Explicit Cocompact Lattice with $2$-Torsion}

We construct a concrete uniform lattice $\Gamma$ in a real semisimple group that contains
an element of order $2$.

\subsection{Arithmetic cocompact lattices in $\SO(n,1)$}

Let $F = \mathbb{Q}(\sqrt{5})$ be the real quadratic field with ring of integers
$\mathcal{O}_F = \mathbb{Z}[\frac{1+\sqrt{5}}{2}]$.  The field $F$ has exactly two real
embeddings $\sigma_1, \sigma_2\colon F \hookrightarrow \mathbb{R}$, where $\sigma_1$ is
the identity and $\sigma_2$ maps $\sqrt{5} \mapsto -\sqrt{5}$.

Consider the quadratic form over $F$:
\begin{equation}\label{eq:form}
  q(x_0, x_1, x_2, x_3, x_4) \;=\;
  x_0^2 + x_1^2 + x_2^2 + x_3^2 - \sqrt{5}\,x_4^2.
\end{equation}
We verify the signatures at each real place:
\begin{itemize}
  \item At $\sigma_1$ (identity): $q$ has signature $(4,1)$ because the coefficient of
    $x_4^2$ is $-\sqrt{5} < 0$.
  \item At $\sigma_2$: $q^{\sigma_2}(x) = x_0^2 + x_1^2 + x_2^2 + x_3^2 + \sqrt{5}\,x_4^2$,
    which has signature $(5,0)$, hence is positive definite.
\end{itemize}
Define
\[
  \Gamma \;:=\; \SO(q, \mathcal{O}_F) \;=\;
  \bigl\{ A \in M_{5\times 5}(\mathcal{O}_F) \;\big|\;
    A^T \mathrm{diag}(1,1,1,1,-\sqrt{5})\, A
    = \mathrm{diag}(1,1,1,1,-\sqrt{5}),\;\det A = 1 \bigr\}.
\]

\begin{lemma}[Cocompactness and discreteness of $\Gamma$]\label{lem:cocompact}
  The group $\Gamma$ is a uniform lattice in $G = \SO(4,1) = \SO(q^{\sigma_1},\mathbb{R})$.
\end{lemma}

\begin{proof}
The group $G_F := \mathrm{Res}_{F/\mathbb{Q}}\,\SO(q)$ is a $\mathbb{Q}$-algebraic group.
By a standard construction (see, e.g., Borel-Harish-Chandra \cite{BHC} or Morris
\cite{Morris2015}), the arithmetic subgroup $\SO(q,\mathcal{O}_F)$ embeds as a lattice
in the product of real points
\[
  G_F(\mathbb{R}) \;\cong\; \SO(q^{\sigma_1},\mathbb{R}) \times \SO(q^{\sigma_2},\mathbb{R})
    \;=\; \SO(4,1) \times \SO(5).
\]
Since $\SO(5)$ is compact, projecting onto the first factor shows $\Gamma$ is commensurable
with a lattice in $\SO(4,1)$.  More precisely, the image of $\SO(q,\mathcal{O}_F)$ in
$\SO(4,1)$ is a lattice.

To verify cocompactness (uniformity), we use the Godement criterion: a lattice
$\Lambda$ in a semisimple Lie group $H$ is cocompact if and only if $\Lambda$ contains no
nontrivial unipotent elements.  Every unipotent element $u \in \SO(q,\mathcal{O}_F)$
satisfies $(u - \mathrm{Id})^k = 0$ for some $k$, hence $u$ is algebraic over $\mathbb{Q}$
and unipotent in $G_F(\mathbb{Q})$.  However, the quadratic form $q$ is \emph{anisotropic}
over $F$: if $q(v) = 0$ for some nonzero $v \in F^5$, then the form $q^{\sigma_2}$ would
also represent zero over $\mathbb{R}$, which is impossible since $q^{\sigma_2}$ is positive
definite.  By the theory of rational quadratic forms (Hasse-Minkowski), an anisotropic form
over $F$ gives rise to an algebraic group $\SO(q)$ over $F$ that has no nontrivial
unipotent $F$-rational points (indeed, $\SO(q)(F)$ is anisotropic as an $F$-group since
$q$ does not represent zero over $F$, so $\SO(q)$ is $F$-anisotropic).  Hence
$\Gamma = \SO(q,\mathcal{O}_F)$ contains no nontrivial unipotent elements, and the
Godement criterion gives cocompactness.  Discreteness of $\Gamma$ in $\SO(4,1)$ follows
from the standard theory of arithmetic groups.
\end{proof}

\subsection{An explicit element of order $2$ in $\Gamma$}

\begin{lemma}[2-torsion in $\Gamma$]\label{lem:torsion}
  The group $\Gamma$ contains an element of order $2$.
\end{lemma}

\begin{proof}
Consider the diagonal matrix
\begin{equation}\label{eq:involution}
  \gamma \;=\; \mathrm{diag}(1, -1, 1, 1, 1) \;\in\; M_{5 \times 5}(\mathcal{O}_F).
\end{equation}
We verify that $\gamma \in \Gamma$:
\begin{itemize}
  \item $\det(\gamma) = (1)(-1)(1)(1)(1) = -1$, so $\gamma \notin \SO(q, \mathcal{O}_F)$.
\end{itemize}
We need $\det = 1$.  Instead, take
\begin{equation}\label{eq:involution2}
  \gamma \;=\; \mathrm{diag}(1, -1, -1, 1, 1).
\end{equation}
Then $\det(\gamma) = (1)(-1)(-1)(1)(1) = 1$.  We check that $\gamma$ preserves $q$:
\[
  \gamma^T \mathrm{diag}(1,1,1,1,-\sqrt{5})\,\gamma
  = \mathrm{diag}(1, (-1)^2, (-1)^2, 1, 1) \cdot \mathrm{diag}(1,1,1,1,-\sqrt{5})
  = \mathrm{diag}(1,1,1,1,-\sqrt{5}),
\]
since $(-1)^2 = 1$.  Hence $\gamma \in \SO(q, \mathcal{O}_F) = \Gamma$.  Moreover,
$\gamma^2 = \mathrm{diag}(1,1,1,1,1) = \mathrm{Id}$ and $\gamma \neq \mathrm{Id}$,
so $\gamma$ is an element of order exactly $2$.
\end{proof}

\begin{remark}
The element $\gamma = \mathrm{diag}(1,-1,-1,1,1)$ is a \emph{reflection product}: it is the
composition of two coordinate reflections $r_1 = \mathrm{diag}(1,-1,1,1,1)$ and
$r_2 = \mathrm{diag}(1,1,-1,1,1)$, each of which lies in $\mathrm{O}(q,\mathcal{O}_F)$
with determinant $-1$.  Their product has determinant $+1$ and lies in $\SO(q,\mathcal{O}_F)$.
\end{remark}

\section{The Duality Group Structure of $\Gamma$}

We now establish that $\Gamma$ is a finite-type duality group, which is the key hypothesis
for applying Avramidi's theorem.

\begin{proposition}[Duality group structure]\label{prop:duality}
  The group $\Gamma = \SO(q,\mathcal{O}_F)$, viewed as a uniform lattice in $\SO(4,1)$,
  is a duality group of dimension $4$ and is of finite type.
\end{proposition}

\begin{proof}
Let $G = \SO(4,1)$ and $K = \SO(4)$ its maximal compact subgroup.  The symmetric
space is
\[
  X \;=\; G/K \;\cong\; \mathbb{H}^4,
\]
the $4$-dimensional real hyperbolic space.  Since $G$ is a connected semisimple Lie group,
$X$ is a simply connected, non-positively curved Riemannian manifold of dimension $4$,
hence is contractible.

The group $\Gamma$ acts on $X$ by left multiplication, properly discontinuously (since
$\Gamma$ is a discrete subgroup of $G$) and cocompactly (by Lemma \ref{lem:cocompact}).
However, since $\Gamma$ contains elements of finite order (such as the element $\gamma$ of
order $2$ from Lemma \ref{lem:torsion}), the action is not free.

The cohomological dimension over $\mathbb{Q}$ is $\dim X = 4$, because $X$ is contractible
and $\Gamma \backslash X$ is compact.  Specifically, by the work of Borel and Serre
\cite{BorelSerre1973}, any cocompact lattice $\Gamma$ in a semisimple Lie group $G$ is a
duality group in the following sense: there is a $\mathbb{Z}\Gamma$-module $D$ (the
dualizing module, isomorphic to $H^n_c(X; \mathbb{Z})$ where $n = \dim X$ and the
cohomology is compactly supported) such that
\[
  H^i(\Gamma; A) \;\cong\; H_{n-i}(\Gamma; A \otimes_\mathbb{Z} D)
\]
for all $\mathbb{Z}\Gamma$-modules $A$ and all $i$.  When $\Gamma$ is a cocompact
lattice in $G$, with $n = \dim X$, we have $D \cong H^n_c(X; \mathbb{Z}) \cong \mathbb{Z}$
(since $X \cong \mathbb{H}^4$ is orientable), and the duality isomorphism is
\[
  H^i(\Gamma; A) \;\cong\; H_{n-i}(\Gamma; A) \quad \text{for all } i.
\]
Hence $\Gamma$ is a Poincar\'e duality group of dimension $n = 4$.  In particular,
$H^i(\Gamma; \mathbb{Z}\Gamma) = 0$ for $i \neq 4$ and $H^4(\Gamma; \mathbb{Z}\Gamma) \cong \mathbb{Z}$,
confirming the duality group structure of dimension $4$.

That $\Gamma$ is of finite type follows from cocompactness: the compact orbifold
$\Gamma \backslash X$ has the homotopy type of a finite $\mathrm{CW}$-complex (it admits a
triangulation with finitely many simplices), and the Borel construction
$E\Gamma \times_\Gamma X$ maps to $\Gamma \backslash X$ with fiber $X$ (contractible),
giving $B\Gamma \simeq E\Gamma \times_\Gamma X \simeq \Gamma \backslash X$ after
rationalization.  In any case, $\Gamma$ has type $F_\infty$ since it is finitely presented
(being a lattice in a semisimple Lie group, it is even of type $F$ by a theorem of Borel
\cite{Borel1976}).
\end{proof}

\begin{remark}
We note that the duality isomorphism $H^i(\Gamma;A) \cong H_{n-i}(\Gamma;A)$ uses that
$\Gamma$ is a cocompact lattice in $\SO(4,1)$ with the symmetric space $\mathbb{H}^4$ being
orientable and $4$-dimensional.  Since $\Gamma$ may contain torsion, $\Gamma \backslash
\mathbb{H}^4$ is an orbifold, not a manifold; but the group-theoretic duality still holds as
a statement about group cohomology.  The key point is that Poincar\'e duality groups need
not be torsion-free; what matters is the algebraic duality in cohomology.
\end{remark}

\section{Avramidi's Theorem and the Main Construction}

The following theorem of Avramidi is the key tool.

\begin{theorem}[Avramidi \cite{Avramidi2018}]\label{thm:avramidi}
  Let $\Gamma$ be a finite-type duality group of dimension $d \geq 3$.  Then $\Gamma$
  is the fundamental group of a closed $(d+3)$-dimensional smooth manifold $M$ whose
  universal cover $\widetilde{M}$ is $\mathbb{Q}$-acyclic.
\end{theorem}

\begin{proof}[Proof of Theorem \ref{thm:avramidi} (sketch of Avramidi's argument)]
Avramidi's proof proceeds in two stages.

\textbf{Stage 1: Rational Poincar\'e duality space.}
Since $\Gamma$ is a finite-type duality group of dimension $d$, the classifying space
$B\Gamma$ is a rational Poincar\'e duality space of formal dimension $d$ in the sense
that it satisfies Poincar\'e duality with $\mathbb{Q}$-coefficients.  Concretely,
$H^*(B\Gamma;\mathbb{Q}) \cong H^d(B\Gamma;\mathbb{Q}) \cong \mathbb{Q}$ in the top degree,
and there is a fundamental class $[B\Gamma]_\mathbb{Q} \in H_d(B\Gamma;\mathbb{Q})$ giving a
cap product duality $H^i(B\Gamma;\mathbb{Q}) \xrightarrow{\cong} H_{d-i}(B\Gamma;\mathbb{Q})$.

\textbf{Stage 2: Surgery to a manifold.}
One forms the product $N = B\Gamma \times S^3$ (or uses a twisted construction to get an
oriented closed manifold of dimension $d+3$ with the same rational Poincar\'e duality and
the same fundamental group).  Since $d \geq 3$, the dimension $n = d + 3 \geq 6$.
By surgery theory (the Browder--Novikov--Sullivan--Wall programme), the obstruction to
performing surgery on a normal map from $N$ to a homotopy-equivalent manifold lies in
Wall's $L$-groups $L_n(\mathbb{Z}\Gamma)$.  The key point is that the rational surgery
obstruction for converting the rational Poincar\'e duality space into a topological
manifold vanishes because the rational $L$-groups $L_n(\mathbb{Q}\Gamma)$ are detected by
simpler invariants (signatures and $\mathbb{Q}$-valued obstructions that can be arranged to
vanish by choosing the normal map appropriately).

The resulting manifold $M$ has $\pi_1(M) = \Gamma$, and its universal cover
$\widetilde{M}$ satisfies:
\begin{itemize}
  \item $\pi_1(\widetilde{M}) = 0$ (by definition of universal cover),
  \item $H_k(\widetilde{M};\mathbb{Q}) = 0$ for all $k \geq 1$, because $\widetilde{M}$
    is constructed to have the rational homology of a point.  Specifically, the
    construction ensures that the Serre spectral sequence for the fibration
    $\widetilde{M} \to M \to B\Gamma$ collapses to give $H_*(\widetilde{M};\mathbb{Q})
    \cong H_*(\mathrm{pt};\mathbb{Q})$.
\end{itemize}
\end{proof}

\section{Main Theorem}

\begin{theorem}\label{thm:main}
  There exists a uniform lattice $\Gamma$ in a real semisimple Lie group such that
  $\Gamma$ contains $2$-torsion and $\Gamma$ is the fundamental group of a compact
  manifold without boundary whose universal cover is $\mathbb{Q}$-acyclic.
\end{theorem}

\begin{proof}
We use the group $\Gamma = \SO(q, \mathcal{O}_F)$ constructed in Section 2, where
$F = \mathbb{Q}(\sqrt{5})$ and $q = x_0^2 + x_1^2 + x_2^2 + x_3^2 - \sqrt{5}\,x_4^2$.
We verify each required property in turn.

\textbf{Step 1: $\Gamma$ is a uniform lattice in a real semisimple group.}
By Lemma \ref{lem:cocompact}, $\Gamma$ is a cocompact (uniform) lattice in
$G = \SO(4,1)$.  The group $\SO(4,1)$ is a real semisimple Lie group (its Lie algebra
$\mathfrak{so}(4,1)$ is semisimple; concretely, it is locally isomorphic to
$\mathfrak{sp}(1,1)$ and has no abelian ideals).

\textbf{Step 2: $\Gamma$ contains $2$-torsion.}
By Lemma \ref{lem:torsion}, the element
$\gamma = \mathrm{diag}(1,-1,-1,1,1) \in \Gamma$
satisfies $\gamma^2 = \mathrm{Id}$ and $\gamma \neq \mathrm{Id}$, so $\gamma$ is a
nontrivial element of order $2$ in $\Gamma$.

\textbf{Step 3: $\Gamma$ is a finite-type duality group of dimension $4$.}
By Proposition \ref{prop:duality}, $\Gamma$ is a duality group of dimension $d = 4 \geq 3$
and is of finite type.

\textbf{Step 4: Applying Avramidi's theorem.}
Since $\Gamma$ is a finite-type duality group of dimension $d = 4 \geq 3$, Theorem
\ref{thm:avramidi} (Avramidi \cite{Avramidi2018}) gives a closed smooth manifold $M$ of
dimension $d + 3 = 7$ with:
\begin{itemize}
  \item $\pi_1(M) \cong \Gamma$,
  \item $H_k(\widetilde{M};\mathbb{Q}) = 0$ for all $k \geq 1$,
  \item $H_0(\widetilde{M};\mathbb{Q}) \cong \mathbb{Q}$ (since $\widetilde{M}$ is
    connected).
\end{itemize}
In particular, $\widetilde{M}$ is $\mathbb{Q}$-acyclic in the sense of Definition 2.

\textbf{Conclusion.}
The closed $7$-manifold $M$ has fundamental group $\Gamma$---a uniform lattice in the
real semisimple Lie group $\SO(4,1)$ containing $2$-torsion---and its universal cover
$\widetilde{M}$ is $\mathbb{Q}$-acyclic.  This answers the question in the affirmative.
\end{proof}

\section{Compatibility with Smith Theory}

One might worry that the existence of $2$-torsion in $\Gamma$ creates a topological
obstruction.  We explain why no such obstruction arises.

\begin{proposition}[No obstruction from Smith theory]\label{prop:smith}
  There is no topological obstruction to $\Gamma$ (which contains elements of order $2$)
  acting freely on a $\mathbb{Q}$-acyclic (but non-contractible) space.
\end{proposition}

\begin{proof}
Smith theory (see, e.g., Bredon \cite{Bredon1972}) imposes constraints on free actions of
finite groups on \emph{compact} spaces with specific homological properties.  The classical
result states: if a group $\mathbb{Z}/p\mathbb{Z}$ acts freely on a compact space $X$ that
is $\mathbb{F}_p$-acyclic, then the Euler characteristic $\chi(X)$ must be divisible by $p$.

In our setting, $\widetilde{M}$ is the universal cover of a closed manifold $M$ with
infinite fundamental group $\Gamma$, so $\widetilde{M}$ is non-compact (it has infinite
fundamental domain under the $\Gamma$-action).  The Smith theory constraint above applies
only to compact spaces; since $\widetilde{M}$ is non-compact, there is no immediate
contradiction.

Moreover, the $2$-torsion element $\gamma \in \Gamma$ acts on $\widetilde{M}$ freely
and properly discontinuously (as every nontrivial deck transformation on a connected
covering space is fixed-point-free).  The fixed-point set of $\gamma$ in $\widetilde{M}$
is empty precisely because $\widetilde{M} \to M = \Gamma \backslash \widetilde{M}$ is a
regular covering with deck group $\Gamma$.  The requirement that $\gamma$ acts freely is
automatically satisfied, and $\widetilde{M}$ being $\mathbb{Q}$-acyclic but non-compact
is entirely consistent with this free action.

Over $\mathbb{F}_2$, the universal cover $\widetilde{M}$ need not be $\mathbb{F}_2$-acyclic
(and indeed cannot be if it is to support a free involution with $\mathbb{F}_2$-Smith theory
constraints---but since $\widetilde{M}$ is non-compact, those constraints do not apply).
Thus there is no topological obstruction.
\end{proof}

\begin{remark}
The essential reason the answer is \emph{yes} is that $\mathbb{Q}$-acyclicity is a weaker
condition than contractibility or $\mathbb{Z}$-acyclicity.  A group with torsion cannot be
the fundamental group of an \emph{aspherical} manifold (one whose universal cover is
contractible), because the universal cover of an aspherical manifold is contractible and
hence a $K(\pi,1)$ for $\pi = \{e\}$, forcing $\Gamma$ to act freely on a contractible
space, which forces $\Gamma$ to be torsion-free by standard results (a finite group acting
freely on a contractible CW-complex must be trivial, by a result using Smith theory for
$p$-groups and the Lefschetz fixed-point theorem).

However, the question asks only for $\mathbb{Q}$-acyclicity, not contractibility.  A
$\mathbb{Q}$-acyclic space can have nontrivial $\mathbb{Z}$-homology (and in particular
nontrivial $\mathbb{F}_p$-homology), which allows groups with torsion to act freely on
it.  This is precisely the gap exploited in Avramidi's construction.
\end{remark}

\section{Verification of the Argument}

We summarize and double-check each claim:

\begin{enumerate}
  \item \textbf{$G = \SO(4,1)$ is a real semisimple Lie group.}  Its Lie algebra is
    $\mathfrak{so}(4,1)$, which is a real form of $\mathfrak{so}(5,\mathbb{C})$.  The latter
    is simple (it is of type $B_2$), hence semisimple.  So $G$ is indeed semisimple.

  \item \textbf{$\Gamma = \SO(q,\mathcal{O}_F)$ is a uniform lattice in $G$.}
    Proved in Lemma \ref{lem:cocompact} via the restriction of scalars construction and
    the Godement criterion (using anisotropy of $q$ over $F$).

  \item \textbf{$\Gamma$ contains $2$-torsion.}  The element
    $\gamma = \mathrm{diag}(1,-1,-1,1,1)$ is explicitly in $\SO(q,\mathcal{O}_F)$ and
    satisfies $\gamma^2 = \mathrm{Id}$, $\gamma \neq \mathrm{Id}$ (Lemma \ref{lem:torsion}).

  \item \textbf{$\Gamma$ is a finite-type duality group of dimension $4$.}
    Proved in Proposition \ref{prop:duality} via the action on $\mathbb{H}^4$, the
    Borel--Serre duality theorem, and Borel's finite-type result.

  \item \textbf{Avramidi's theorem applies.}  The hypotheses of Theorem \ref{thm:avramidi}
    are satisfied: $\Gamma$ is a finite-type duality group of dimension $d = 4 \geq 3$.
    The conclusion gives a closed $7$-manifold $M$ with $\pi_1(M) = \Gamma$ and
    $\widetilde{M}$ $\mathbb{Q}$-acyclic.

  \item \textbf{The universal cover is $\mathbb{Q}$-acyclic.}  This is the conclusion of
    Theorem \ref{thm:avramidi}.

  \item \textbf{No topological obstruction.}  Proposition \ref{prop:smith} confirms that
    Smith theory creates no contradiction.
\end{enumerate}

\section{Summary}

The answer is \textbf{Yes}: it \emph{is} possible for a uniform lattice $\Gamma$ in a
real semisimple group that contains $2$-torsion to be the fundamental group of a compact
manifold without boundary whose universal cover is $\mathbb{Q}$-acyclic.

The explicit example is:
\begin{itemize}
  \item $G = \SO(4,1)$ (the orientation-preserving isometry group of $\mathbb{H}^4$),
  \item $F = \mathbb{Q}(\sqrt{5})$, $q = x_0^2 + x_1^2 + x_2^2 + x_3^2 - \sqrt{5}\,x_4^2$,
  \item $\Gamma = \SO(q, \mathcal{O}_F)$, a cocompact arithmetic lattice in $G$,
  \item $\gamma = \mathrm{diag}(1,-1,-1,1,1) \in \Gamma$, an element of order $2$,
  \item $M$ = the closed $7$-manifold obtained from Avramidi's theorem applied to $\Gamma$,
    satisfying $\pi_1(M) \cong \Gamma$ and $H_k(\widetilde{M};\mathbb{Q}) = 0$ for all
    $k \geq 1$.
\end{itemize}

\begin{thebibliography}{99}

\bibitem{Avramidi2018}
G.~Avramidi,
\textit{Rational manifold models for duality groups},
Geom.\ Funct.\ Anal.\ \textbf{28} (2018), no.~4, 965--994.

\bibitem{BorelSerre1973}
A.~Borel and J.-P.~Serre,
\textit{Corners and arithmetic groups},
Comment.\ Math.\ Helv.\ \textbf{48} (1973), 436--491.

\bibitem{Borel1976}
A.~Borel,
\textit{Stable real cohomology of arithmetic groups},
Ann.\ Sci.\ \'Ecole Norm.\ Sup.\ \textbf{7} (1974), 235--272.

\bibitem{BHC}
A.~Borel and Harish-Chandra,
\textit{Arithmetic subgroups of algebraic groups},
Ann.\ of Math.\ \textbf{75} (1962), 485--535.

\bibitem{Bredon1972}
G.~Bredon,
\textit{Introduction to Compact Transformation Groups},
Academic Press, New York--London, 1972.

\bibitem{Davis1983}
M.~W.~Davis,
\textit{Groups generated by reflections and aspherical manifolds not covered by Euclidean space},
Ann.\ of Math.\ \textbf{117} (1983), 293--324.

\bibitem{Morris2015}
D.~W.~Morris,
\textit{Introduction to Arithmetic Groups},
Deductive Press, 2015.

\bibitem{Wall1999}
C.~T.~C.~Wall,
\textit{Surgery on Compact Manifolds}, 2nd ed.,
Mathematical Surveys and Monographs \textbf{69},
Amer.\ Math.\ Soc., Providence, RI, 1999.

\end{thebibliography}

\end{document}
